\documentclass[12pt]{article}
\usepackage{amsmath, amssymb}
\usepackage{geometry}
\geometry{margin=1in}

\title{CSE400: Lecture 17 Scribe\\
Joint Random Variables \& Joint Distributions}
\author{}
\date{}

\begin{document}

\maketitle

\section{Why Study Joint Random Variables?}

In science and real-life applications, we are often interested in two (or more) random variables simultaneously. \text{\cite{ppt}} :contentReference[oaicite:0]{index=0}

\subsection*{Examples}
\begin{itemize}
    \item Height and weight of giraffes
    \item IQ and birthweight of children
    \item Frequency of exercise and rate of heart disease
    \item Air pollution level and respiratory illness rate
    \item Number of Facebook friends and age of members
\end{itemize}

\subsection*{Applications}
\begin{itemize}
    \item Smart Transportation: $X=$ vehicle speed, $Y=$ traffic density
    \item Wireless Networks: $X=$ SNR, $Y=$ throughput
    \item Weather Prediction: $X=$ rainfall, $Y=$ temperature
    \item Drone Delivery: $X=$ wind speed, $Y=$ delivery time
    \item Dice Example: conditions such as $X+Y=7$, $X>Y$
\end{itemize}

Joint distributions allow us to compute probabilities involving both variables and understand their relationship. If variables are not independent, dependence is studied using covariance and correlation.

\section{Discrete Case: Joint PMF}

\subsection{Setup}
Let $X$ and $Y$ be discrete random variables:
\[
X = \{x_1, x_2, \dots, x_n\}, \quad Y = \{y_1, y_2, \dots, y_m\}
\]

Then:
\[
(X,Y) \in \{(x_i, y_j)\}
\]

\subsection{Definition}
\[
p(x_i, y_j) = P(X = x_i, Y = y_j)
\]

\subsection{Properties}
\begin{itemize}
    \item $0 \leq p(x_i, y_j) \leq 1$
    \item $\sum_{i=1}^{n} \sum_{j=1}^{m} p(x_i, y_j) = 1$
\end{itemize}

\subsection{Joint PMF Table}
Joint PMF is represented in tabular form where rows correspond to values of $X$ and columns correspond to values of $Y$.

\section{Example: Rolling Two Dice}

\subsection{Experiment}
Roll two fair dice.

\[
X = \text{value on first die}, \quad T = \text{sum of both dice}
\]

Each outcome $(i,j)$ has probability:
\[
\frac{1}{36}
\]

\subsection{Joint PMF Construction}
The joint PMF table of $(X,T)$ contains values:
\[
p(x,t) =
\begin{cases}
\frac{1}{36}, & \text{if valid combination} \\
0, & \text{otherwise}
\end{cases}
\]

\subsection{Observation}
\begin{itemize}
    \item Entries are either $0$ or $\frac{1}{36}$
    \item $X$ and $T$ are not independent
\end{itemize}

\section{Continuous Case: Joint PDF}

\subsection{Analogy}
\begin{itemize}
    \item Discrete values $\rightarrow$ Continuous intervals
    \item PMF $\rightarrow$ PDF
    \item Summation $\rightarrow$ Double integration
\end{itemize}

\subsection{Domain}
If:
\[
X \in [a,b], \quad Y \in [c,d]
\]
then:
\[
(X,Y) \in [a,b] \times [c,d]
\]

\subsection{Definition}
A function $f(x,y)$ is a joint PDF if:
\[
P((X,Y) \text{ lies in a small rectangle}) \approx f(x,y)\,dx\,dy
\]

\subsection{Properties}
\begin{itemize}
    \item $f(x,y) \geq 0$
    \item $\int_c^d \int_a^b f(x,y)\,dx\,dy = 1$
\end{itemize}

\textbf{Note:} $f(x,y)$ may be greater than 1 since it is a density.

\section{Worked Example}

Given:
\[
f(x,y) = 4xy, \quad 0 \leq x \leq 1, \quad 0 \leq y \leq 1
\]

\subsection{Step 1: Check Validity}
\[
\int_0^1 \int_0^1 4xy \, dx\, dy
= \int_0^1 \left( \int_0^1 4xy\,dx \right) dy
\]

\[
= \int_0^1 4y \left( \int_0^1 x\,dx \right) dy
= \int_0^1 4y \cdot \frac{1}{2} \, dy
= \int_0^1 2y\,dy
\]

\[
= \left[ y^2 \right]_0^1 = 1
\]

Hence, $f(x,y)$ is a valid joint PDF.

\subsection{Step 2: Compute $P(A)$}

\[
A = \{X < 0.5, Y > 0.5\}
\]

\[
P(A) = \int_0^{0.5} \int_{0.5}^{1} 4xy\, dy\, dx
\]

First integrate w.r.t $y$:
\[
\int_{0.5}^{1} 4xy\,dy = 4x \left[\frac{y^2}{2}\right]_{0.5}^{1}
= 4x \cdot \frac{1 - 0.25}{2}
= 1.5x
\]

Now integrate w.r.t $x$:
\[
\int_0^{0.5} 1.5x\,dx = 1.5 \cdot \frac{x^2}{2} \Big|_0^{0.5}
= \frac{3}{16}
\]

\[
\boxed{P(A) = \frac{3}{16}}
\]

\section{Joint Cumulative Distribution Function (CDF)}

\subsection{Definition}
\[
F(x,y) = P(X \leq x, Y \leq y)
\]

\section{Joint CDF: Continuous Case}

\[
F(x,y) = \int_c^y \int_a^x f(u,v)\,du\,dv
\]

\subsection{Recovering PDF}
\[
f(x,y) = \frac{\partial^2}{\partial x \partial y} F(x,y)
\]

\section{Joint CDF: Discrete Case}

\[
F(x,y) = \sum_{x_i \leq x} \sum_{y_j \leq y} p(x_i,y_j)
\]

\subsection{Interpretation}
\begin{itemize}
    \item Add probabilities from rows with $x_i \leq x$
    \item Add probabilities from columns with $y_j \leq y$
\end{itemize}

\section{Properties of Joint CDF}

\begin{itemize}
    \item $F(x,y)$ is non-decreasing
    \item $\lim_{(x,y)\to(-\infty,-\infty)} F(x,y) = 0$
    \item $\lim_{(x,y)\to(\infty,\infty)} F(x,y) = 1$
\end{itemize}

\end{document}